% =====================================================================
%  lecon.tex — compile UNE leçon du recueil, seule.
%
%  À placer à la racine du recueil (à côté de « Oraux Agreg ext maths.tex ») :
%  il charge le même préambule (Preambule/preambule.tex) et \input le même
%  fichier de leçon que le recueil, donc rien n'est dupliqué.
%
%  Usage : depuis la racine du recueil,  ./lecon.sh 101
%  Le script fait les deux passes de pdflatex (nombre de pages) et écrit le PDF
%  dans Lecons/AG/101.pdf (algèbre, leçons 1xx) ou Lecons/AP/201.pdf
%  (analyse-probas, leçons 2xx). À la main, c'est l'option -output-directory :
%    pdflatex -output-directory=Lecons/AG -jobname=101 "\def\numlecon{101}% =====================================================================
%  lecon.tex — compile UNE leçon du recueil, seule.
%
%  À placer à la racine du recueil (à côté de « Oraux Agreg ext maths.tex ») :
%  il charge le même préambule (Preambule/preambule.tex) et \input le même
%  fichier de leçon que le recueil, donc rien n'est dupliqué.
%
%  Usage : depuis la racine du recueil,  ./lecon.sh 101
%  Le script fait les deux passes de pdflatex (nombre de pages) et écrit le PDF
%  dans Lecons/AG/101.pdf (algèbre, leçons 1xx) ou Lecons/AP/201.pdf
%  (analyse-probas, leçons 2xx). À la main, c'est l'option -output-directory :
%    pdflatex -output-directory=Lecons/AG -jobname=101 "\def\numlecon{101}% =====================================================================
%  lecon.tex — compile UNE leçon du recueil, seule.
%
%  À placer à la racine du recueil (à côté de « Oraux Agreg ext maths.tex ») :
%  il charge le même préambule (Preambule/preambule.tex) et \input le même
%  fichier de leçon que le recueil, donc rien n'est dupliqué.
%
%  Usage : depuis la racine du recueil,  ./lecon.sh 101
%  Le script fait les deux passes de pdflatex (nombre de pages) et écrit le PDF
%  dans Lecons/AG/101.pdf (algèbre, leçons 1xx) ou Lecons/AP/201.pdf
%  (analyse-probas, leçons 2xx). À la main, c'est l'option -output-directory :
%    pdflatex -output-directory=Lecons/AG -jobname=101 "\def\numlecon{101}% =====================================================================
%  lecon.tex — compile UNE leçon du recueil, seule.
%
%  À placer à la racine du recueil (à côté de « Oraux Agreg ext maths.tex ») :
%  il charge le même préambule (Preambule/preambule.tex) et \input le même
%  fichier de leçon que le recueil, donc rien n'est dupliqué.
%
%  Usage : depuis la racine du recueil,  ./lecon.sh 101
%  Le script fait les deux passes de pdflatex (nombre de pages) et écrit le PDF
%  dans Lecons/AG/101.pdf (algèbre, leçons 1xx) ou Lecons/AP/201.pdf
%  (analyse-probas, leçons 2xx). À la main, c'est l'option -output-directory :
%    pdflatex -output-directory=Lecons/AG -jobname=101 "\def\numlecon{101}\input{lecon}"
%  Sans rien passer en ligne de commande, « pdflatex lecon » compile la leçon
%  par défaut indiquée ci-dessous, dans le dossier courant.
%
%  Le fichier de la leçon est lu au même endroit que dans le recueil, dans le
%  dossier des centaines : 101 -> 100/101, 201 -> 200/201. Pour un autre
%  chemin (un développement par exemple) :  ./lecon.sh DevAlg/baseburnside
%  (le PDF est alors écrit à côté du fichier source).
% =====================================================================
\documentclass[a4paper,10pt,oneside,twocolumn,landscape]{book}

% ---------- Leçon à compiler ----------
\ifdefined\fichier
	% chemin explicite fourni en ligne de commande : rien à faire
\else
	\providecommand{\numlecon}{101}   % leçon par défaut (ignorée si \numlecon vient de la ligne de commande)
	\edef\fichier{\the\numexpr(\numlecon-50)/100*100\relax/\numlecon}
\fi
\edef\titrepdf{\ifdefined\numlecon Leçon \numlecon\else\fichier\fi}

\input{Preambule/preambule}

% ---------- Ajustements propres à une leçon compilée seule ----------
% Pas de chapitre ici : la section (titre de la leçon) n'est pas numérotée (sinon « 0.1 ») et les sous-sections sont numérotées 1, 2, 3...
\titleformat{\section}{\normalfont\Large\bfseries}{}{0pt}{}
\renewcommand{\thesubsection}{\arabic{subsection}}
% Pas de page de titre ici : numérotation des pages sans le décalage de +1 du recueil.
\fancyfoot[C]{Page \thepage/\pageref{LastPage}}
\hypersetup{pdftitle={\titrepdf}}

\begin{document}

\input{\fichier}

\end{document}
"
%  Sans rien passer en ligne de commande, « pdflatex lecon » compile la leçon
%  par défaut indiquée ci-dessous, dans le dossier courant.
%
%  Le fichier de la leçon est lu au même endroit que dans le recueil, dans le
%  dossier des centaines : 101 -> 100/101, 201 -> 200/201. Pour un autre
%  chemin (un développement par exemple) :  ./lecon.sh DevAlg/baseburnside
%  (le PDF est alors écrit à côté du fichier source).
% =====================================================================
\documentclass[a4paper,10pt,oneside,twocolumn,landscape]{book}

% ---------- Leçon à compiler ----------
\ifdefined\fichier
	% chemin explicite fourni en ligne de commande : rien à faire
\else
	\providecommand{\numlecon}{101}   % leçon par défaut (ignorée si \numlecon vient de la ligne de commande)
	\edef\fichier{\the\numexpr(\numlecon-50)/100*100\relax/\numlecon}
\fi
\edef\titrepdf{\ifdefined\numlecon Leçon \numlecon\else\fichier\fi}

\input{Preambule/preambule}

% ---------- Ajustements propres à une leçon compilée seule ----------
% Pas de chapitre ici : la section (titre de la leçon) n'est pas numérotée (sinon « 0.1 ») et les sous-sections sont numérotées 1, 2, 3...
\titleformat{\section}{\normalfont\Large\bfseries}{}{0pt}{}
\renewcommand{\thesubsection}{\arabic{subsection}}
% Pas de page de titre ici : numérotation des pages sans le décalage de +1 du recueil.
\fancyfoot[C]{Page \thepage/\pageref{LastPage}}
\hypersetup{pdftitle={\titrepdf}}

\begin{document}

\input{\fichier}

\end{document}
"
%  Sans rien passer en ligne de commande, « pdflatex lecon » compile la leçon
%  par défaut indiquée ci-dessous, dans le dossier courant.
%
%  Le fichier de la leçon est lu au même endroit que dans le recueil, dans le
%  dossier des centaines : 101 -> 100/101, 201 -> 200/201. Pour un autre
%  chemin (un développement par exemple) :  ./lecon.sh DevAlg/baseburnside
%  (le PDF est alors écrit à côté du fichier source).
% =====================================================================
\documentclass[a4paper,10pt,oneside,twocolumn,landscape]{book}

% ---------- Leçon à compiler ----------
\ifdefined\fichier
	% chemin explicite fourni en ligne de commande : rien à faire
\else
	\providecommand{\numlecon}{101}   % leçon par défaut (ignorée si \numlecon vient de la ligne de commande)
	\edef\fichier{\the\numexpr(\numlecon-50)/100*100\relax/\numlecon}
\fi
\edef\titrepdf{\ifdefined\numlecon Leçon \numlecon\else\fichier\fi}

\input{Preambule/preambule}

% ---------- Ajustements propres à une leçon compilée seule ----------
% Pas de chapitre ici : la section (titre de la leçon) n'est pas numérotée (sinon « 0.1 ») et les sous-sections sont numérotées 1, 2, 3...
\titleformat{\section}{\normalfont\Large\bfseries}{}{0pt}{}
\renewcommand{\thesubsection}{\arabic{subsection}}
% Pas de page de titre ici : numérotation des pages sans le décalage de +1 du recueil.
\fancyfoot[C]{Page \thepage/\pageref{LastPage}}
\hypersetup{pdftitle={\titrepdf}}

\begin{document}

\input{\fichier}

\end{document}
"
%  Sans rien passer en ligne de commande, « pdflatex lecon » compile la leçon
%  par défaut indiquée ci-dessous, dans le dossier courant.
%
%  Le fichier de la leçon est lu au même endroit que dans le recueil, dans le
%  dossier des centaines : 101 -> 100/101, 201 -> 200/201. Pour un autre
%  chemin (un développement par exemple) :  ./lecon.sh DevAlg/baseburnside
%  (le PDF est alors écrit à côté du fichier source).
% =====================================================================
\documentclass[a4paper,10pt,oneside,twocolumn,landscape]{book}

% ---------- Leçon à compiler ----------
\ifdefined\fichier
	% chemin explicite fourni en ligne de commande : rien à faire
\else
	\providecommand{\numlecon}{101}   % leçon par défaut (ignorée si \numlecon vient de la ligne de commande)
	\edef\fichier{\the\numexpr(\numlecon-50)/100*100\relax/\numlecon}
\fi
\edef\titrepdf{\ifdefined\numlecon Leçon \numlecon\else\fichier\fi}

\input{Preambule/preambule}

% ---------- Ajustements propres à une leçon compilée seule ----------
% Pas de chapitre ici : la section (titre de la leçon) n'est pas numérotée (sinon « 0.1 ») et les sous-sections sont numérotées 1, 2, 3...
\titleformat{\section}{\normalfont\Large\bfseries}{}{0pt}{}
\renewcommand{\thesubsection}{\arabic{subsection}}
% Pas de page de titre ici : numérotation des pages sans le décalage de +1 du recueil.
\fancyfoot[C]{Page \thepage/\pageref{LastPage}}
\hypersetup{pdftitle={\titrepdf}}

\begin{document}

\input{\fichier}

\end{document}
